\documentclass[11pt,a4paper]{article}
\usepackage[margin=1in]{geometry}
\usepackage[T1]{fontenc}
\usepackage[utf8]{inputenc}
\usepackage{lmodern}
\usepackage{microtype}
\usepackage{amsmath,amssymb,mathtools}
\usepackage[hidelinks]{hyperref}
\usepackage{fancyhdr}
\setlength{\headheight}{14pt}
\pagestyle{fancy}
\fancyhf{}
\lhead{Ernst Steinitz}
\rhead{Deutscher Text}
\cfoot{\thepage}
\setlength{\parindent}{1.5em}
\setlength{\parskip}{0.35em}
\allowdisplaybreaks
\newcommand{\Cl}{\operatorname{Cl}}
\newcommand{\Md}{\operatorname{Md}}
\newcommand{\Nm}{\operatorname{Nm}}
\newcommand{\Rg}{\operatorname{Rg}}
\newcommand{\Gd}{\operatorname{Gd}}
\newcommand{\Dg}{\operatorname{Dg}}
\newcommand{\tht}{\vartheta}
\begin{document}
\section*{Zur Theorie der Moduln}
\begin{center}
Von \emph{Ernst Steinitz} in Charlottenburg.
\end{center}
\section*{§ 4. Isomorphismus; die Function $\chi$.}

\subsection*{Definition des Isomorphismus.}

Lässt sich zwischen den Resten $\rho$ eines Moduls $\mathfrak a$ und den Resten $\rho'$ eines Moduls $\mathfrak a'$ eine solche eindeutige Beziehung herstellen, dass, wenn $\rho_i,\rho_i'$ und $\rho_k,\rho_k'$ zwei Paare zugeordneter Reste sind, stets $\rho_i+\rho_k$, $\rho_i'+\rho_k'$ ebenfalls ein Paar zugeordneter Reste bilden, so wird die getroffene Zuordnung eine ,,isomorphe`` Beziehung zwischen $\mathfrak a$ und $\mathfrak a'$ genannt. Zwei Moduln $\mathfrak a,\mathfrak a'$, welche wenigstens eine isomorphe Beziehung zulassen, heissen ,,zu einander isomorph``.\footnote{Die Bezeichnung isomorph ist der Gruppentheorie entnommen.} Die Beziehung des Isomorphismus ist also stets eine wechselseitige.

Sind $\rho_1,\rho_2,\ldots$ die sämmtlichen Reste von $\mathfrak a$, $\rho_1',\rho_2',\ldots$ in dieser Reihenfolge die isomorph zugeordneten Reste von $\mathfrak a'$, so werde die vorliegende isomorphe Beziehung durch
\[
 \rho_i\parallel \rho_i'
\]
bezeichnet. Aus der Definition des Isomorphismus folgt leicht, dass allgemein
\[
 \sum_i c_i\rho_i,
 \qquad
 \sum_i c_i\rho_i'
\]
entsprechende Reste von $\mathfrak a$ und $\mathfrak a'$ sind. Ist $\mathfrak a'$ isomorph zu $\mathfrak a''$ und stellt
\[
 \rho_i'\parallel \rho_i''
\]
eine isomorphe Beziehung zwischen $\mathfrak a'$ und $\mathfrak a''$ dar, so ist
\[
 \rho_i\parallel \rho_i''
\]
eine isomorphe Beziehung zwischen $\mathfrak a$ und $\mathfrak a''$, welche ,,aus $\rho_i\parallel\rho_i'$ und $\rho_i'\parallel\rho_i''$ zusammengesetzt`` heisst. Setzt man die isomorphe Beziehung $\rho_i\parallel\rho_i'$ nach einander mit den sämmtlichen verschiedenen isomorphen Beziehungen von $\mathfrak a'$ zu $\mathfrak a''$ zusammen, so erhält man ebensoviele isomorphe Beziehungen zwischen $\mathfrak a$ und $\mathfrak a''$, und mit diesen sind alle erschöpft. Bezeichnet man demnach mit
\[
 \chi(\mathfrak a\parallel\mathfrak a')
\]
die Anzahl der zwischen zwei Moduln $\mathfrak a$ und $\mathfrak a'$ bestehenden isomorphen Beziehungen, so gilt der Satz:

Ist $\mathfrak a$ isomorph zu $\mathfrak a'$, $\mathfrak a'$ isomorph zu $\mathfrak a''$, so ist $\mathfrak a$ isomorph zu $\mathfrak a''$ und
\[
 \chi(\mathfrak a\parallel\mathfrak a'')=
 \chi(\mathfrak a'\parallel\mathfrak a'')=
 \chi(\mathfrak a\parallel\mathfrak a').
\]
Die Moduln $\mathfrak a,\mathfrak a',\mathfrak a''$ brauchen natürlich nicht alle verschieden zu sein. Nimmt man $\mathfrak a=\mathfrak a''$ an, so erhält man, den Isomorphismus von $\mathfrak a$ und $\mathfrak a'$ vorausgesetzt,
\[
 \chi(\mathfrak a\parallel\mathfrak a')=
 \chi(\mathfrak a\parallel\mathfrak a)=
 \chi(\mathfrak a'\parallel\mathfrak a').
\]

Eine isomorphe Beziehung eines Moduls $\mathfrak a$ auf sich selbst kann als eine Permutation unter seinen Resten gedeutet werden. Die so erhaltenen Permutationen bilden eine Gruppe; $\chi(\mathfrak a\parallel\mathfrak a)$, wofür wir auch einfach
\[
 \chi(\mathfrak a)
\]
schreiben, ist die Ordnung dieser Gruppe.

Es soll nun näher untersucht werden, unter welchen Bedingungen zwei Moduln $\mathfrak a$ und $\mathfrak a'$ isomorph sind. Man sieht leicht ein, dass dies der Fall ist, wenn sie derselben Classe angehören. Ist nämlich alsdann $[E]$ eine unimodulare Transformation, welche $\mathfrak a$ in $\mathfrak a'$ überführt, so wird (§ 2) durch dieselbe auch jeder Rest von $\mathfrak a$ in einen Rest von $\mathfrak a'$ verwandelt und auf diese Weise eine eindeutige Beziehung zwischen den Resten von $\mathfrak a$ und $\mathfrak a'$ hergestellt, die man als eine isomorphe sofort erkennt.

Umgekehrt möge jetzt gezeigt werden, dass zwei isomorphe Moduln desselben Grades $\mathfrak a,\mathfrak a'$ auch derselben Classe angehören müssen. Dabei kann auf Grund der vorangegangenen Betrachtungen, unbeschadet der Allgemeinheit des Beweises, vorausgesetzt werden, dass $\mathfrak a$ und $\mathfrak a'$ innerhalb ihrer Classen die Hauptmoduln sind. Es möge ferner der Fall, dass die Normen von $\mathfrak a$ und $\mathfrak a'$ verschwinden, ausgeschlossen werden, da derselbe für unseren Zweck nicht in Betracht kommt. Bezeichnet
\[
 \rho_i\parallel \rho_i'
\]
eine isomorphe Beziehung zwischen $\mathfrak a$ und $\mathfrak a'$, ist $m$ irgend eine ganze Zahl, und setzt man für $\rho_i,\rho_i'$ successive alle Reste von $\mathfrak a,\mathfrak a'$, so gehört das System $m\rho_i$ ebenso oft dem Modul $\mathfrak a$ an, als das System $m\rho_i'$ dem Modul $\mathfrak a'$ angehört. Die Berechnung der gesuchten Anzahl gestaltet sich aber, wenn $\mathfrak a,\mathfrak a'$ als Hauptmoduln vorausgesetzt werden, besonders einfach. Sind $a_1,a_2,\ldots,a_n$ die Invarianten von $\mathfrak a$, $a_1',a_2',\ldots,a_n'$ diejenigen von $\mathfrak a'$, so erkennt man sogleich, dass die Anzahl der Reste $\rho_i$, für welche $m\rho_i$ in $\mathfrak a$ enthalten ist, gleich
\[
 (m,a_1)(m,a_2)\cdots(m,a_n)
\]
wird, und ebenso ergiebt sich für die Anzahl der Reste $\rho_i'$, für welche $m\rho_i'$ dem Modul $\mathfrak a'$ angehört, der Ausdruck
\[
 (m,a_1')(m,a_2')\cdots(m,a_n').
\]
Demnach muss für jede Zahl $m$ die Gleichung
\[
 (m,a_1)(m,a_2)\cdots(m,a_n)=(m,a_1')(m,a_2')\cdots(m,a_n')
\]
bestehen, und dies kann, wie eine einfache Ueberlegung zeigt, nur der Fall sein, wenn
\[
 a_1=a_1',\quad a_2=a_2',\quad\ldots,\quad a_n=a_n'
\]
ist, wenn also $\mathfrak a$ und $\mathfrak a'$ derselben Classe angehören.

Sehr wohl aber können zwei Moduln verschiedenen Grades isomorph sein; die Frage nach dem Isomorphismus zweier Moduln wird durch folgenden Satz erledigt.

\emph{I.} Zwei Moduln sind dann und nur dann isomorph, wenn ihre Invariantensysteme nach Ausscheidung der Invarianten, welche sich auf $1$ reduciren, identisch sind.

Um zu beweisen, dass die ausgesprochene Bedingung eine hinreichende ist, nehmen wir an, es habe der Modul $\mathfrak a$ die Invarianten
\[
 a_1=1,\ldots,a_r=1,a_{r+1},a_{r+2},\ldots,a_n,
\]
der Modul $\mathfrak a'$ die Invarianten
\[
 a_1'=a_{r+1},\quad a_2'=a_{r+2},\quad\ldots,\quad a_{n-r}'=a_n,
\]
dann ist zu zeigen, dass $\mathfrak a$ und $\mathfrak a'$ isomorph sind. Wir dürfen dabei $\mathfrak a$ und $\mathfrak a'$ als Hauptmoduln voraussetzen, und unter dieser Annahme gestaltet sich der Beweis sehr einfach. Bilden nämlich jetzt
\[
 \rho_1,\rho_2,\ldots
\]
ein vollständiges Restsystem in Bezug auf $\mathfrak a$, und wird unter $\rho_i'$ die Zeile verstanden, welche aus $\rho_i$ durch Fortlassung der ersten $r$ Elemente hervorgeht, so ergiebt sich ohne Weiteres, dass
\[
 \rho_1',\rho_2',\ldots
\]
ein vollständiges Restsystem von $\mathfrak a'$ bilden, und dass durch
\[
 \rho_i\parallel\rho_i'
\]
eine isomorphe Beziehung zwischen $\mathfrak a$ und $\mathfrak a'$ dargestellt wird. Verbindet man das soeben erhaltene Resultat mit dem früheren, wonach Moduln gleichen Grades nur dann isomorph sind, wenn sie zu derselben Classe gehören, so wird die Nothwendigkeit der im Satze I ausgesprochenen Bedingung klar.

Für die Function $\chi$ können aus diesem Satze einige einfache Folgerungen gezogen werden. Zunächst zeigt es sich, dass für alle Moduln $\mathfrak a$ einer Classe $A$ $\chi$ denselben Werth hat, sodass das Symbol $\chi(\mathfrak a)$ zweckmässig wird durch
\[
 \chi(A)
\]
ersetzt werden können. Sodann hat man den Satz:

\emph{II.} Wenn die Invariantensysteme der Classen $A$ und $A'$ nach Ausscheidung der sich auf $1$ reducirenden Invarianten identisch sind, so ist
\[
 \chi(A)=\chi(A').
\]

\section*{§ 5. Zusammenhang zwischen den Functionen $\varphi$, $\chi$, $\psi$.}

Die von uns betrachteten Functionen $\varphi,\chi,\psi$ stehen in einem sehr einfachen Zusammenhange, welcher durch die folgende Betrachtung klar wird.

Ist $A$ eine Classe $n$ten Grades, $\mathfrak a$ ein Modul von $A$, $R$ ein $r$-zeiliger Rest von $\mathfrak a$, so besteht, wie eine einfache Ueberlegung zeigt, der Satz: Damit
\begin{equation}
 (\mathfrak a,R)=\mathfrak e
\tag{1}
\end{equation}
sei, ist nothwendig und hinreichend, dass man zu jeder Zeile $\rho$ von $n$ Elementen eine Zeile $\sigma$ von $r$ Elementen finden könne, welche der Congruenz
\begin{equation}
 \sigma R\equiv\rho\pmod{\mathfrak a}
\tag{2}
\end{equation}
genügt. Nehmen wir nun an, die Gleichung (1) sei erfüllt und es sei $r=n$, dann bilden die sämmtlichen $\sigma$, für welche $\sigma R$ in $\mathfrak a$ enthalten ist, einen Modul $\mathfrak a'$ vom Grade $n$; die sämmtlichen $\sigma$, welche einer Congruenz von der Form (2), worin $\rho$ als gegeben zu betrachten ist, genügen, sind Repräsentanten eines bestimmten Restes von $\mathfrak a'$, und umgekehrt: für alle Repräsentanten eines bestimmten Restes von $\mathfrak a'$ stellt $\sigma R$ einen und denselben Rest von $\mathfrak a$ dar. Man erhält daher eine eindeutige Beziehung zwischen den Resten von $\mathfrak a$ und $\mathfrak a'$, indem man jedem Rest $\sigma$ von $\mathfrak a'$ den Rest $\sigma R$ von $\mathfrak a$ entsprechen lässt, und diese Beziehung ist, wie man leicht sieht, eine isomorphe. Daraus folgt, dass der Modul $\mathfrak a'$ der Classe $A$ angehört (§ 4, I). Wir wollen sagen, der Rest $R$ von $\mathfrak a$ gehöre zum Modul $\mathfrak a'$ und wollen die durch $R$ vermittelte isomorphe Beziehung mit $[R]$ bezeichnen.

Ist $R'$ ebenfalls ein $n$-zeiliger der Bedingung (1) genügender Rest von $\mathfrak a$, aber von $R$ verschieden, so kann $R'$ wohl zu demselben Modul $\mathfrak a'$ gehören, aber die isomorphe Beziehung $[R']$ ist von $[R]$ verschieden. Anderenfalls nämlich müsste für jede Zeile $\sigma$ die Congruenz
\[
 \sigma R\equiv\sigma R'\pmod{\mathfrak a}
\]
bestehen, und hieraus würde man, indem man für $\sigma$ nach einander die Zeilen des Systems
\[
 E^0=\begin{pmatrix}
 1&0&0&\cdots&0\\
 0&1&0&\cdots&0\\
 0&0&1&\cdots&0\\
 \vdots&\vdots&\vdots&&\vdots\\
 0&0&0&\cdots&1
 \end{pmatrix}
\]
setzt, die Congruenz
\[
 R'\equiv R\pmod{\mathfrak a}
\]
folgern.

Andererseits lässt sich leicht zeigen, dass, wenn man für $R$ jeden $n$-zeiligen der Bedingung (1) genügenden Rest von $\mathfrak a$ setzt, durch das Symbol $[R]$ jede zwischen $\mathfrak a$ und irgend einem Modul $\mathfrak a'$ von $A$ mögliche isomorphe Beziehung einmal dargestellt wird. Ist nämlich eine solche Beziehung zwischen $\mathfrak a$ und $\mathfrak a'$ gegeben, werden durch dieselbe den durch die $n$ Zeilen von $E^0$ repräsentirten Resten von $\mathfrak a'$ die Reste $\rho_1,\rho_2,\ldots,\rho_n$ von $\mathfrak a$ zugeordnet und ist $R$ der von diesen gebildete $n$-zeilige Rest, so ergiebt die Definition des Isomorphismus, dass dem Rest $\sigma$ von $\mathfrak a'$ der Rest $\sigma R$ von $\mathfrak a$ entsprechen muss. Da somit jeder Rest von $\mathfrak a$ in der Form $\sigma R$ darstellbar ist, so ergiebt sich zunächst, dass $R$ der Bedingung
\[
 (\mathfrak a,R)=\mathfrak e
\]
genügt, und sodann weiter, dass die vorliegende isomorphe Beziehung mit der durch das Symbol $[R]$ bezeichneten identisch ist.

Das Resultat dieser Betrachtungen ist, dass jeder der Reste $R$, deren Anzahl gleich $\varphi(A)$ ist, zu einem bestimmten der $\psi(A)$ Moduln der Classe $A$ gehört, und dass zu jedem Modul $\mathfrak a'$ von $A$ $\chi(A)$ Reste $R$ gehören, so viele nämlich, als isomorphe Beziehungen zwischen $\mathfrak a$ und $\mathfrak a'$ existiren. So gelangt man zu der wichtigen Relation
\[
\text{I.}\qquad \varphi(A)=\chi(A)\cdot\psi(A).
\]
Aus dieser und den Formeln (§ 3, II; § 2, XVII) erhält man noch die Beziehung
\[
 \chi(A)=\prod_p \chi(A_p),
\]
welche sich natürlich auch leicht direct nachweisen liesse.

\section*{§ 6. Bestimmung der Functionen $\psi$ und $\chi$.}

Auf Grund der Sätze § 2, IX; § 4, II; § 5, I ist es nun leicht, die Functionen $\psi$ und $\chi$ zunächst als Quotienten aus Producten von $\varphi$-Functionen darzustellen.

Ist $A=\Cl(a_1\mid a_2\mid\cdots\mid a_n)$ irgend eine Classe $n$ten Grades von positiver Norm, und setzt man
\[
 A'=\Cl\left(1\mid {a_2\over a_1}\mid {a_3\over a_1}\mid\cdots\mid {a_n\over a_1}\right),
\]
\[
 A_i=\Cl\left({a_{i+1}\over a_i}\mid {a_{i+2}\over a_i}\mid\cdots\mid {a_n\over a_i}\right),
\qquad
 A_i'=\Cl\left(1\mid {a_{i+2}\over a_{i+1}}\mid {a_{i+3}\over a_{i+1}}\mid\cdots\mid {a_n\over a_{i+1}}\right),
\]
sodass $A_i$, $A_i'$ Classen vom Grade $n-i$ sind, so wird
\[
 A=a_1A',\qquad A_i={a_{i+1}\over a_i}A_i',
\]
\[
 \psi(A)=\psi(A'),\quad \psi(A_i)=\psi(A_i')\quad (§ 2,\ IX),
\]
\[
 \psi(A')={\varphi(A')\over\chi(A')},\quad
 \psi(A_i')={\varphi(A_i')\over\chi(A_i')}\quad (§ 5,\ I),
\]
\[
 \chi(A')=\chi(A_1),\quad \chi(A_i')=\chi(A_{i+1})\quad (§ 4,\ II),
\]
\[
 \chi(A_1)={\varphi(A_1)\over\psi(A_1)},\quad
 \chi(A_{i+1})={\varphi(A_{i+1})\over\psi(A_{i+1})}\quad (§ 5,\ I),
\]
mithin
\[
 \psi(A)={\varphi(A')\over\varphi(A_1)}\psi(A_1),
 \qquad
 \psi(A_i)={\varphi(A_i')\over\varphi(A_{i+1})}\psi(A_{i+1}).
\]
Da $\psi(A_{n-1})=1$ ist, so folgt aus diesen Relationen
\[
\text{I.}\qquad
 \psi(A)={\varphi(A')\varphi(A_1')\cdots\varphi(A_{n-2}')
 \over
 \varphi(A_1)\varphi(A_2)\cdots\varphi(A_{n-1})},
\]
\[
\text{II.}\qquad
 \chi(A)={\varphi(A)\varphi(A_1)\cdots\varphi(A_{n-1})
 \over
 \varphi(A')\varphi(A_1')\cdots\varphi(A_{n-2}')}.
\]

Die weitere Ausrechnung, welche wir nur für die Function $\psi$ durchführen wollen, erfordert, zunächst den Fall der primären Classe zu erledigen. Es sei also
\[
 \Cl(a_1\mid a_2\mid\cdots\mid a_n)=A=A_p=
 \Cl(p^{r_1}\mid p^{r_2}\mid\cdots\mid p^{r_n})
\]
und somit
\[
 r_1\le r_2\le\cdots\le r_n.
\]
Mit $e_1,e_2,\ldots,e_k$ seien die verschiedenen unter den Exponenten $r$ in ihrer natürlichen Reihenfolge bezeichnet. Kommt $e_1$ $s_1$mal, $e_2$ $s_2$mal, \ldots, $e_k$ $s_k$mal unter den Exponenten $r$ vor, so ist
\begin{equation}
 s_1+s_2+\cdots+s_k=n.
\tag{1}
\end{equation}
Unter den Differenzen $r_m-r_l$, in denen $m>l$ ist, ist die Differenz $e_i-e_h$ $(i>h)$ $s_i\cdot s_h$mal enthalten. Somit ist
\begin{equation}
 \sum_{i>h}(e_i-e_h)s_i s_h=\sum_{m>l}r_m-r_l
 =(1-n)r_1+(3-n)r_2+(5-n)r_3+\cdots+(n-3)r_{n-1}+(n-1)r_n.
\tag{2}
\end{equation}
Setzt man nun
\[
 s_1+s_2+\cdots+s_i=t_i\qquad (i=1,2,\ldots,k),
\]
so folgt aus
\begin{equation}
 p^{e_1}=a_1=a_2=\cdots=a_{t_1},\quad
 p^{e_2}=a_{t_1+1}=\cdots=a_{t_2},\quad\ldots,\quad
 p^{e_k}=a_{t_{k-1}+1}=\cdots=a_{t_k},
\tag{3}
\end{equation}
dass in (I) nur die Classen
\[
 A',\ A_1',\ldots,A_{k-2}'
\quad\text{und}\quad
 A_1,\ldots,A_{k-1}
\]
wesentlich verschieden auftreten. Genauer erhält man
\begin{equation}
 \psi(A)=
 {\varphi(A')\varphi(A_1')\cdots\varphi(A_{k-2}')
 \over
 \varphi(A_1)\varphi(A_2)\cdots\varphi(A_{k-1})}.
\tag{4}
\end{equation}
Es haben aber die Classen $A'$, $A_i$, $A_i'$ beziehungsweise die Grade $n$, $n-t_i$, $n-t_i$, während die Anzahlen ihrer durch $p$ theilbaren Invarianten beziehungsweise gleich $n-t_1$, $n-t_i$, $n-t_{i+1}$ sind; folglich ist (§ 3, IIIa)
\[
 \varphi(A')=(\Nm A')^n {\tht(p;n)\over\tht(p;t_1)}=(\Nm A')^n {\tht(p;n)\over\tht(p;s_1)},
\]
\[
 \varphi(A_i)=(\Nm A_i)^{n-t_i}\tht(p;n-t_i),
\]
\[
 \varphi(A_i')=(\Nm A_i')^{n-t_i}
 {\tht(p;n-t_i)\over\tht(p;t_{i+1}-t_i)}
 = (\Nm A_i')^{n-t_i}{\tht(p;n-t_i)\over\tht(p;s_{i+1})}.
\]
Also nach (4)
\begin{equation}
 \psi(A)=
 { (\Nm A')^n(\Nm A_1')^{n-t_1}\cdots(\Nm A_{k-2}')^{n-t_{k-2}}
 \over
 (\Nm A_1)^{n-t_1}(\Nm A_2)^{n-t_2}\cdots(\Nm A_{k-1})^{n-t_{k-1}} }
 {\tht(p;n)\over\tht(p;s_1)\tht(p;s_2)\cdots\tht(p;s_k)}.
\tag{5}
\end{equation}
Nun ist
\[
 { (\Nm A_{i-1}')^{n-t_{i-1}}
 \over (\Nm A_i)^{n-t_i} }
 =\left({\Nm A_{i-1}'\over\Nm A_i'}\right)^{n-t_i}(\Nm A_{i-1}')^{t_i-t_{i-1}},
\]
daher
\begin{equation}
 { (\Nm A')^n(\Nm A_1')^{n-t_1}\cdots(\Nm A_{k-2}')^{n-t_{k-2}}
 \over
 (\Nm A_1)^{n-t_1}(\Nm A_2)^{n-t_2}\cdots(\Nm A_{k-1})^{n-t_{k-1}} }
 =\left({\Nm A'\over\Nm A_1}\right)^{t_1}
 \left({\Nm A_1'\over\Nm A_2}\right)^{t_2}
 \cdots
 \left({\Nm A_{k-3}'\over\Nm A_{k-2}}\right)^{t_{k-2}}
 (\Nm A_{k-2}')^{t_{k-1}}.
\tag{6}
\end{equation}
Ferner ergiebt sich
\[
 {\Nm A_{i-1}'\over\Nm A_i'}=\left({a_{t_i+1}\over a_{t_i}}\right)^{n-t_i},
 \qquad
 \Nm A_{k-2}'=\left({a_n\over a_{t_{k-1}}}\right)^{n-t_{k-1}},
\]
\begin{equation}
 \prod_{i=1}^{k-2}\left({\Nm A_{i-1}'\over\Nm A_i'}\right)^{t_i}
 (\Nm A_{k-2}')^{t_{k-1}}
 =\prod_{i=1}^{k-1}\left({a_{t_i+1}\over a_{t_i}}\right)^{(n-t_i)t_i}
 =\prod_{i=1}^{n-1}\left({a_{i+1}\over a_i}\right)^{i(n-i)}
 =\prod_{i=1}^n a_i^{2i-1-n}.
\tag{7}
\end{equation}
Aus (5), (6), (7) erhält man jetzt
\[
\text{III.}\qquad
 \psi(A)=a_1^{1-n}a_2^{3-n}a_3^{5-n}\cdots a_{n-1}^{n-3}a_n^{n-1}
 {\tht(p;n)\over\tht(p;s_1)\tht(p;s_2)\cdots\tht(p;s_k)}.
\]
oder, unter Berücksichtigung von (2), (3),
\[
\text{IIIa.}\qquad
 \psi(A)=p^{\sum_{i>h}s_i s_h(e_i-e_h)}
 {\tht(p;n)\over\tht(p;s_1)\tht(p;s_2)\cdots\tht(p;s_k)}.
\]

Führt man eine Function $g(x;n,m)$\footnote{Es ist dieselbe, welche in der Gauss'schen Abhandlung ,,Summatio quarundam serierum singularium`` mit $(n,m)$ bezeichnet wird.} ein durch die Gleichung
\begin{equation}
 g(x;n,m)=
 { (1-x^n)(1-x^{n-1})\cdots(1-x^{n-m+1})
 \over
 (1-x)(1-x^2)\cdots(1-x^m) },
\tag{8}
\end{equation}
worin $n$ eine ganze, $m$ eine positive ganze Zahl bedeutet, so ist
\begin{equation}
 g(x;n,m)=g(x;n-1,m)+x^{n-m}g(x;n-1,m-1),
\tag{9}
\end{equation}
und diese Gleichung gilt auch für beliebiges ganzzahliges $m$, wenn festgesetzt wird, dass
\begin{equation}
 g(x;n,0)=1,
\tag{10}
\end{equation}
\begin{equation}
 g(x;n,m)=0\qquad (m<0)
\tag{11}
\end{equation}
sein soll. Es gilt ferner die Beziehung
\begin{equation}
 g(x;n,m)=x^{m(n-m)}g\left({1\over x};n,m\right).
\tag{12}
\end{equation}
Aus den Gleichungen (9), (10), (11) folgt, dass $g(x;n,m)$ für $n\ge0$ eine ganze Function von $x$ ist, welche nur ganzzahlige nichtnegative Coefficienten besitzt. Für $n\ge m\ge0$ besteht zwischen den Functionen $g$ und $\tht$ die Relation
\begin{equation}
 g\left({1\over x};n,m\right)={\tht(x;n)\over\tht(x;m)\tht(x;n-m)}.
\tag{13}
\end{equation}
Wird zur Abkürzung
\begin{equation}
 \psi(A)=p^{\sum_{i>h}s_i s_h(e_i-e_h)}\bar\psi(A)
 =a_1^{1-n}a_2^{3-n}\cdots a_n^{n-1}\bar\psi(A)
\tag{14}
\end{equation}
gesetzt, so ergeben sich für die Functionen $\psi$, $\bar\psi$ die folgenden Ausdrücke
\begin{equation}
 \bar\psi(A)=
 g\left({1\over p};n,s_1\right)
 g\left({1\over p};n-s_1,s_2\right)
 \cdots
 g\left({1\over p};n-s_1-\cdots-s_{k-2},s_{k-1}\right),
\tag{15}
\end{equation}
\begin{equation}
\text{IIIb.}\qquad
 \psi(A)=p^{\sum_{i>h}s_i s_h(e_i-e_h-1)}
 g(p;n,s_1)g(p;n-s_1,s_2)\cdots
 g(p;n-s_1-\cdots-s_{k-2},s_{k-1}).
\tag{16}
\end{equation}
Aus dieser Darstellung erkennt man, dass $\psi(A)$ eine ganze Function von $p$ ist, ferner, dass der Coefficient der höchsten Potenz von $p$ gleich 1 ist und dass die übrigen Coefficienten ganze nichtnegative Zahlen sind.

Hat man eine beliebige Classe $A$ von positiver Norm, so wird durch § 2, XVII die Berechnung von $\psi(A)$ auf den Fall der primären Classe zurückgeführt. Setzt man
\[
 \prod_p\bar\psi(A_p)=\bar\psi(A),
\]
so ergiebt sich
\[
 \psi(A)=\prod_p\left((a_1)_p^{1-n}(a_2)_p^{3-n}\cdots(a_n)_p^{n-1}\bar\psi(A_p)\right)
 =a_1^{1-n}a_2^{3-n}\cdots a_n^{n-1}\bar\psi(A),
\]
sodass Gleichung (14) allgemeine Geltung besitzt.

Wendet man die gewonnenen Resultate auf den Fall $n=2$ an, so erhält man für $A=\Cl(a_1\mid a_2)$
\[
 \psi(A)={a_2\over a_1}\bar\psi(A)
 ={a_2\over a_1}\prod_{p\mid a_2/a_1}\bar\psi(A_p)
 ={a_2\over a_1}\prod_{p\mid a_2/a_1}{\tht(p;2)\over(\tht(p;1))^2}
 ={a_2\over a_1}\prod_{p\mid a_2/a_1}\left(1+{1\over p}\right),
\]
wo das Product sich auf die in $a_2/a_1$ aufgehenden Primzahlen erstreckt.

\emph{Beispiel:} Für
\[
 a_1=3,\\ a_2=90
\]
ist
\[
 {a_2\over a_1}=30=2\cdot3\cdot5,
\]
\[
 \psi(A)=30\cdot {3\over2}\cdot {4\over3}\cdot {6\over5}=72.
\]
Für
\[
 n=3,
 \qquad A=\Cl(a_1\mid a_2\mid a_3)
\]
ist
\[
 \psi(A)=\left({a_3\over a_1}\right)^2\bar\psi(A)
 =\left({a_3\over a_1}\right)^2
 \prod_{p\mid a_3/a_1}{\tht(p;3)\over\tht(p;1)\tht(p;2)}
 \prod_{\substack{p\mid a_2/a_1\\ p\mid a_3/a_2}}
 {\tht(p;2)\over\tht(p;1)\tht(p;1)},
\]
also
\[
 \psi(A)=\left({a_3\over a_1}\right)^2
 \prod_{p\mid a_3/a_1}\left(1+{1\over p}+{1\over p^2}\right)
 \prod_{\substack{p\mid a_2/a_1\\ p\mid a_3/a_2}}
 \left(1+{1\over p}\right),
\]
wo das erste Product sich auf die in $a_3/a_1$, das zweite auf die in $a_2/a_1$ und $a_3/a_2$ zugleich aufgehenden Primzahlen erstreckt.

\emph{Beispiel:} Für
\[
 a_1=2,\qquad a_2=210,\qquad a_3=6930
\]
ist
\[
 {a_3\over a_1}=3465=3^2\cdot5\cdot7\cdot11,
 \qquad {a_2\over a_1}=105=3\cdot5\cdot7,
 \qquad {a_3\over a_2}=33=3\cdot11.
\]
\[
 \psi(A)=(3465)^2\cdot {13\over9}\cdot {31\over25}\cdot {57\over49}\cdot {133\over121}\cdot {4\over3}=36661716.
\]

\section*{§ 7. Einführung einiger Functionen, welche von mehreren Classen abhängen.}

Die Function $\psi$ bildet die Grundlage für die Bestimmung anderer Functionen, welche für die Theorie der Moduln von Wichtigkeit sind und welche im Folgenden behandelt werden sollen. Dabei wollen wir nur Moduln und Classen von positiver Norm in Betracht ziehen.

Ein Modul $\mathfrak a$ von positiver Norm hat, wie man leicht erkennt, nur eine endliche Anzahl von Divisoren, welche sich (§ 2, III) auf diejenigen Classen vertheilen, welche Divisoren von $\Cl\mathfrak a$ sind. Wir wollen nun diejenigen Divisoren von $\mathfrak a$ ins Auge fassen, welche einer bestimmten Classe $B$ angehören. Sind
\[
 \mathfrak b_1,\mathfrak b_2,\ldots,\mathfrak b_k
\]
diese Moduln und bezeichnet $E$ ein unimodulares System, so sind (§ 2, VI)
\[
 \mathfrak b_1E,\mathfrak b_2E,\ldots,\mathfrak b_kE
\]
die sämmtlichen Divisoren von $\mathfrak aE$ innerhalb der Classe $B$. Hieraus folgt, dass die Anzahl der Divisoren von $\mathfrak a$ innerhalb $B$ nur von den Classen $A=\Cl\mathfrak a$ und $B$ abhängt. Wir bezeichnen diese Anzahl mit
\[
 t(A\mid B).
\]
In gleicher Weise erkennt man, dass die Anzahl der Vielfachen, welche ein Modul $\mathfrak b$ der Classe $B$ innerhalb der Classe $A$ besitzt, allein eine Function der Classen $A$ und $B$ ist. Sie sei mit
\[
 v(A\mid B)
\]
bezeichnet. Die Anzahl der Modulpaare $\mathfrak a_i,\mathfrak b_k$, welche den Bedingungen
\[
 \Cl\mathfrak a_i=A,
 \qquad
 \Cl\mathfrak b_k=B,
 \qquad
 (\mathfrak a_i,\mathfrak b_k)=\mathfrak b_k
\]
genügen, wird durch jeden der Ausdrücke $\psi(A)t(A\mid B)$, $\psi(B)v(A\mid B)$ dargestellt, also ist
\[
\text{I.}\qquad
 \psi(A)t(A\mid B)=\psi(B)v(A\mid B).
\]
Ferner ergeben sich aus § 2, XV die Formeln
\[
\text{II.}\qquad t(A\mid B)=\prod_p t(A_p\mid B_p),
\]
\[
\text{III.}\qquad v(A\mid B)=\prod_p v(A_p\mid B_p).
\]
Nach § 2, III wissen wir, dass $t(A\mid B)$ und $v(A\mid B)$ von 0 verschieden oder gleich 0 sind, je nachdem $A$ durch $B$ theilbar ist oder nicht. Endlich gelten für jede Classe $A$ die Beziehungen (§ 2, IV)
\[
\text{IV.}\qquad t(A\mid A)=v(A\mid A)=1.
\]

Zu anderen Classenfunctionen wird man durch die folgenden Betrachtungen geführt.

Sind die Moduln $\mathfrak a_1,\mathfrak a_2,\ldots,\mathfrak a_s$, welche nicht alle verschieden zu sein brauchen, Divisoren eines Moduls $\mathfrak c$, so ist auch ihr kleinstes gemeinsames Vielfaches $[\mathfrak a_1,\mathfrak a_2,\ldots,\mathfrak a_s]$ ein Divisor von $\mathfrak c$. Man kann sich nun die Aufgabe stellen, zu gegebenem $\mathfrak c$ $s$ Moduln $\mathfrak a_1,\mathfrak a_2,\ldots,\mathfrak a_s$ so zu bestimmen, dass
\[
 [\mathfrak a_1,\mathfrak a_2,\ldots,\mathfrak a_s]=\mathfrak c
\]
wird, eine Aufgabe, die natürlich immer und im Allgemeinen auf vielfache Weise lösbar ist. Etwas Anderes aber ist es, wenn man die Bedingung stellt, dass jeder der Moduln $\mathfrak a_i$ einer bestimmten gegebenen Classe $A_i$ angehören soll. Aus § 2, VII folgt der Satz: Ist es für irgend einen Modul $\mathfrak c$ der Classe $\Gamma$ möglich, ihn als kleinstes gemeinsames Vielfaches von $s$ Moduln $\mathfrak a_1,\mathfrak a_2,\ldots,\mathfrak a_s$ darzustellen, welche beziehungsweise den Classen $A_1,A_2,\ldots,A_s$ angehören, so ist dasselbe für jeden Modul der Classe $\Gamma$ möglich. Allgemeiner: Die Anzahl der Systeme von je $s$ Moduln $\mathfrak a_1,\mathfrak a_2,\ldots,\mathfrak a_s$, welche den Bedingungen
\[
 \Cl\mathfrak a_1=A_1,
 \quad
 \Cl\mathfrak a_2=A_2,
 \quad\ldots\quad
 \Cl\mathfrak a_s=A_s;
 \qquad
 [\mathfrak a_1,\mathfrak a_2,\ldots,\mathfrak a_s]=\mathfrak c
\]
genügen, ist für jeden Modul $\mathfrak c$ der Classe $\Gamma$ die nämliche. Wir bezeichnen diese Anzahl mit
\[
 \bar\gamma(\Gamma;A_1,A_2,\ldots,A_s),
\]
sie stellt eine Function der Classen $\Gamma,A_1,\ldots,A_s$ dar, die in Bezug auf $A_1,\ldots,A_s$ symmetrisch ist. Die Frage, ob ein Modul $\mathfrak c$ der Classe $\Gamma$ sich darstellen lasse als kleinstes gemeinsames Vielfaches von $s$ Moduln $\mathfrak a_1,\ldots,\mathfrak a_s$, welche beziehungsweise den Classen $A_1,\ldots,A_s$ angehören, fällt hiernach zusammen mit der Frage, ob $\bar\gamma(\Gamma;A_1,\ldots,A_s)>0$ ist. Die Entscheidung hierüber lässt sich natürlich in jedem gegebenen Falle durch eine endliche Anzahl von Versuchen herbeiführen. Es handelt sich aber darum, ob man nicht allgemeine, leicht zu übersehende Beziehungen aufstellen könne, deren Bestehen zwischen den Classen $\Gamma,A_1,\ldots,A_s$ die nothwendige und hinreichende Bedingung für das Nichtverschwinden der Function $\bar\gamma(\Gamma;A_1,\ldots,A_s)$ darstellt.

Nun erkennt man zunächst, dass $\bar\gamma(\Gamma;A_1,\ldots,A_s)$ nur dann von 0 verschieden sein kann, wenn $\Gamma$ durch jede der Classen $A_1,\ldots,A_s$ theilbar ist. Betrachtet man daher die Classen $A_1,\ldots,A_s$ als gegeben, so stellt ihr kleinstes gemeinsames Vielfaches $\Gamma'=[A_1,\ldots,A_s]$ gewissermassen das Minimum unter denjenigen Classen $\Gamma$ dar, für welche $\bar\gamma(\Gamma;A_1,\ldots,A_s)>0$ ausfällt, indem nämlich erstens $\bar\gamma(\Gamma';A_1,\ldots,A_s)>0$ ist, weil der Hauptmodul von $\Gamma'$ das kleinste gemeinsame Vielfache der Hauptmoduln von $A_1,\ldots,A_s$ ist, und zweitens jede der Bedingung $\bar\gamma(\Gamma;A_1,\ldots,A_s)>0$ genügende Classe $\Gamma$ ein Vielfaches von $\Gamma'$ ist. Eine weitere Untersuchung zeigt, dass unter den in Rede stehenden Classen $\Gamma$ auch ein Maximum existirt, d. h. eine unter ihnen, sie heisse $\Gamma''$, ist durch alle theilbar. Die Bildungsweise der Classe $\Gamma''$ soll bald angegeben werden.

Zu ganz analogen Resultaten wird man bei den Untersuchungen über den grössten gemeinsamen Theiler von Moduln geführt. Wir bezeichnen mit
\[
 \bar\delta(\Delta;A_1,\ldots,A_s)
\]
die Anzahl der Systeme von je $s$ Moduln $\mathfrak a_1,\ldots,\mathfrak a_s$, welche beziehungsweise den Classen $A_1,\ldots,A_s$ entnommen sind und einen gegebenen Modul $\mathfrak b$ der Classe $\Delta$ zum grössten gemeinsamen Theiler haben. Betrachtet man die Classen $A_1,\ldots,A_s$ als gegeben, so liegen die sämmtlichen Classen $\Delta$, für welche $\bar\delta(\Delta;A_1,\ldots,A_s)>0$ ist, zwischen einem Maximum $\Delta'$ und einem Minimum $\Delta''$, sodass
\[
 \bar\delta(\Delta';A_1,\ldots,A_s)>0,
 \qquad
 \bar\delta(\Delta'';A_1,\ldots,A_s)>0,
\]
und jede andere Classe der betrachteten Art Divisor von $\Delta'$ und Vielfaches von $\Delta''$ ist. Hierbei ist $\Delta'$ offenbar der grösste gemeinsame Theiler der Classen $A_1,\ldots,A_s$.

Um zu den Classen $\Gamma''$, $\Delta''$ zu gelangen, hat man folgendermassen zu verfahren. Es sei
\[
 A_i=\Cl(a_{i1}\mid a_{i2}\mid\cdots\mid a_{in})\qquad (i=1,2,\ldots,s).
\]
Man bilde aus den sämmtlichen $s\cdot n$ Invarianten
\[
 a_{11},\ldots,a_{1n},\ldots,a_{s1},\ldots,a_{sn}
\]
der Classen $A_1,\ldots,A_s$ ein Diagonalsystem $K$ und bestimme dessen Classe $K$. Diese ist vom Grade $s\cdot n$ und offenbar unabhängig von der Reihenfolge der Classen $A_1,\ldots,A_s$. Die ersten $n$ Invarianten von $K$ bilden das Invariantensystem einer Classe $n$ten Grades, welche wir den ,,Grenztheiler`` der Classen $A_1,\ldots,A_s$ nennen und mit
\[
 (\!|A_1,\ldots,A_s|\!)
\]
bezeichnen. Ebenso bilden die letzten $n$ Invarianten von $K$ das Invariantensystem einer Classe $n$ten Grades. Dieselbe heisse das ,,Grenzvielfache`` der Classen $A_1,\ldots,A_s$ und sei mit
\[
 [\!|A_1,\ldots,A_s|\!]
\]
bezeichnet. Sie bildet das oben erwähnte Maximum für die Classen $\Gamma$, bei denen $\bar\gamma(\Gamma;A_1,\ldots,A_s)>0$ ist, und ebenso ist der Grenztheiler von $A_1,\ldots,A_s$ das Minimum der Classen $\Delta$, für welche $\bar\delta(\Delta;A_1,\ldots,A_s)>0$ wird. Es ist nicht schwer, den Nachweis hierfür zu geben, ohne auf eine allgemeine Bestimmung der Functionen $\bar\gamma$ und $\bar\delta$ einzugehen. Es ersteht aber jetzt die Frage, ob auch umgekehrt für alle zwischen $[\!|A_1,\ldots,A_s|\!]$ und $[A_1,\ldots,A_s]$ gelegenen Classen $\Gamma$ und für alle zwischen $(A_1,\ldots,A_s)$ und $(\!|A_1,\ldots,A_s|\!)$ gelegenen Classen $\Delta$ die Functionen $\bar\gamma(\Gamma;A_1,\ldots,A_s)$ und $\bar\delta(\Delta;A_1,\ldots,A_s)>0$ ausfallen. Die Entscheidung hierüber ist mir erst durch ein eingehendes Studium der Functionen $\bar\gamma$ und $\bar\delta$ möglich gewesen; es hat sich gezeigt, dass die angeregte Frage zu bejahen ist. Damit hat man alsdann den folgenden Doppelsatz gewonnen, welcher namentlich auch für die Theorie der Gruppen von vertauschbaren Elementen von Interesse ist:

\emph{V. a.} Damit ein Modul $\mathfrak c$ der Classe $\Gamma$ sich darstellen lasse als kleinstes gemeinsames Vielfaches von $s$ Moduln $\mathfrak a_1,\ldots,\mathfrak a_s$, welche beziehungsweise den Classen $A_1,\ldots,A_s$ angehören, ist nothwendig und hinreichend, dass $\Gamma$ zwischen dem Grenzvielfachen und dem kleinsten gemeinsamen Vielfachen der Classen $A_1,\ldots,A_s$ liegt.

\emph{b.} Damit ein Modul $\mathfrak b$ der Classe $\Delta$ sich darstellen lasse als grösster gemeinsamer Theiler von $s$ Moduln $\mathfrak a_1,\ldots,\mathfrak a_s$, welche beziehungsweise den Classen $A_1,\ldots,A_s$ angehören, ist nothwendig und hinreichend, dass $\Delta$ zwischen dem grössten gemeinsamen Theiler und dem Grenztheiler der Classen $A_1,\ldots,A_s$ liegt.

Im Folgenden wird der Nachweis dieses Satzes vollständig durch die Untersuchungen über die Functionen $\bar\gamma$ und $\bar\delta$ erledigt.

Aus § 2, XVI folgen die Relationen
\[
\text{VI.}\qquad
 \bar\gamma(\Gamma;A_1,\ldots,A_s)
 =\prod_p \bar\gamma(\Gamma_p;(A_1)_p,\ldots,(A_s)_p),
\]
\[
 \bar\delta(\Delta;A_1,\ldots,A_s)
 =\prod_p \bar\delta(\Delta_p;(A_1)_p,\ldots,(A_s)_p).
\]
Man ersieht hieraus auch leicht, dass man, um den Nachweis von V zu führen, nur nöthig hat, ihn für primäre Classen zu erbringen.

\section*{§ 8. Die Functionen $t$ und $v$ dargestellt durch die Function $\psi$.}

Wir wenden uns nun einer näheren Betrachtung der Functionen $t(A\mid B)$ und $v(A\mid B)$ zu und wollen zunächst zeigen, wie sich dieselben mit Hülfe von $\psi$-Functionen darstellen lassen. Es kann dies in verschiedener Weise geschehen.

Es sei $A=\Cl(a_1\mid a_2\mid\cdots\mid a_n)$, $B=\Cl(b_1\mid b_2\mid\cdots\mid b_n)$, $\mathfrak a$ ein beliebiger Modul von $A$, endlich seien
\begin{equation}
 \mathfrak b_1,\mathfrak b_2,\ldots,\mathfrak b_{\psi(B)}
\tag{1}
\end{equation}
die sämmtlichen Moduln von $B$. Der Modul $\mathfrak a$ ist durch $a_1\mathfrak e$ theilbar (§ 2, X); ist also $\mathfrak a$ durch $\mathfrak b_k$ theilbar, so ist $\mathfrak a$ auch theilbar durch $[a_1\mathfrak e,\mathfrak b_k]$, und umgekehrt. Die Moduln
\begin{equation}
 [a_1\mathfrak e,\mathfrak b_1],
 [a_1\mathfrak e,\mathfrak b_2],\ldots,
 [a_1\mathfrak e,\mathfrak b_{\psi(B)}]
\tag{2}
\end{equation}
gehören sämmtlich der Classe $[a_1E,B]$ an (§ 2, XI), und eine einfache Ueberlegung zeigt, dass jeder Modul von $[a_1E,B]$ in der Reihe (2) gleich oft, mithin
\[
 {\psi(B)\over\psi([a_1E,B])}
\]
mal vorkommt. Da $\mathfrak a$ durch $t(A\mid [a_1E,B])$ Moduln der Classe $[a_1E,B]$ theilbar ist, so erhält man den Satz:

\emph{I.} Ist $a_1$ die erste Invariante der Classe $A$, so ist
\[
 t(A\mid B)={\psi(B)\over\psi([a_1E,B])}\cdot t(A\mid [a_1E,B]).
\]
Ferner gilt der Satz:

\emph{II.} Haben die Classen
\[
 A=\Cl(a_1\mid a_2\mid\cdots\mid a_n)
 \quad\text{und}\quad
 B=\Cl(b_1\mid b_2\mid\cdots\mid b_n)
\]
dieselbe erste Invariante $a_1=b_1$ und wird
\[
 A'=\Cl(a_2\mid\cdots\mid a_n),
 \qquad
 B'=\Cl(b_2\mid\cdots\mid b_n)
\]
gesetzt, so ist
\[
 t(A\mid B)=t(A'\mid B').
\]

\emph{Beweis:} Es sei $A$ das Hauptsystem von $A$, $\alpha$ seine erste Zeile, $\Md A=\mathfrak a$, $\mathfrak a'$ der Hauptmodul von $A'$, $\mathfrak b=\Md B$ irgend ein Divisor von $\mathfrak a$ innerhalb der Classe $B$. Ist
\[
 B=\begin{pmatrix}
 b_{11}&\cdots&b_{1n}\\
 \vdots&&\vdots\\
 b_{n1}&\cdots&b_{nn}
 \end{pmatrix}
\]
und werden mit $\beta_1,\beta_2,\ldots,\beta_n$ die Zeilen von $B$ bezeichnet, so ist zunächst $\mathfrak b=\Md(\beta_1,\ldots,\beta_n)$ und, weil $\mathfrak a$ durch $\mathfrak b$ theilbar ist, mithin die Zeile $\alpha$ dem Modul $\mathfrak b$ angehört,
\[
 \mathfrak b=\Md(\alpha,\beta_1,\ldots,\beta_n).
\]
Aus $\Cl B=B$, $a_1=b_1$ folgt, dass die Elemente $b_{i1}$ sämmtlich durch $a_1$ theilbar sind. Setzt man nun $b_{i1}=m_i a_1$ $(i=1,\ldots,n)$, $\beta_i-m_i\alpha=\gamma_i$, so wird
\[
 \mathfrak b=\Md(\alpha,\gamma_1,\ldots,\gamma_n).
\]
In $\gamma_i$ ist das erste Element $=0$; bezeichnet $\gamma_i'$ die nach Weglassung dieser Null zurückbleibende Zeile von $n-1$ Elementen, so ist $\mathfrak b'=\Md(\gamma_1',\ldots,\gamma_n')$ ein Modul $(n-1)$ten Grades. Ist
\begin{equation}
 B'=\begin{pmatrix}
 b_{11}'&\cdots&b_{1,n-1}'\\
 \vdots&&\vdots\\
 b_{n-1,1}'&\cdots&b_{n-1,n-1}'
 \end{pmatrix}
\tag{3}
\end{equation}
irgend eine quadratische Basis dieses Moduls, so ist
\begin{equation}
 \bar B=\begin{pmatrix}
 a_1&0&\cdots&0\\
 0&b_{11}'&\cdots&b_{1,n-1}'\\
 \vdots&\vdots&&\vdots\\
 0&b_{n-1,1}'&\cdots&b_{n-1,n-1}'
 \end{pmatrix}
\tag{4}
\end{equation}
eine Basis von $\mathfrak b$. Aus $\Cl\bar B=B$ folgt leicht $\Cl\mathfrak b'=\Cl B'=B'$, und aus der Theilbarkeit von $\mathfrak a$ durch $\mathfrak b$ ergiebt sich, dass $\mathfrak a'$ durch $\mathfrak b'$ theilbar ist. Umgekehrt ist Folgendes leicht zu ersehen: Hat man zwei quadratische Systeme $B'$, $\bar B$, wie sie durch (3) und (4) dargestellt werden, und ist $\Md B'=\mathfrak b'$ Divisor von $\mathfrak a'$, $\Cl B'=B'$, so folgt, dass $\Md\bar B=\mathfrak b$ Divisor von $\mathfrak a$, $\Cl\bar B=B$ ist. Endlich erkennt man, dass in der angegebenen Weise der Modul $\mathfrak b'$ durch den Modul $\mathfrak b$ und ebenso der Modul $\mathfrak b$ durch den Modul $\mathfrak b'$ eindeutig bestimmt ist. Es hat also $\mathfrak a$ in der Classe $B$ ebensoviele Divisoren wie $\mathfrak a'$ in der Classe $B'$, d. h. es ist
\[
 t(A\mid B)=t(A'\mid B').
\]

Es sei jetzt $A=\Cl(a_1\mid a_2\mid\cdots\mid a_n)$ theilbar durch $B=\Cl(b_1\mid b_2\mid\cdots\mid b_n)$ und es werde
\begin{equation}
\left\{
\begin{aligned}
 \Cl(a_{i+1}\mid a_{i+2}\mid\cdots\mid a_n)&=A_i,\\
 \Cl([a_i,b_{i+1}]\mid[a_i,b_{i+2}]\mid\cdots\mid[a_i,b_n])&=B_i,\\
 \Cl([a_{i+1},b_{i+1}]\mid[a_{i+1},b_{i+2}]\mid\cdots\mid[a_{i+1},b_n])&=\bar B_i,
\end{aligned}
\right.
\quad (i=0,1,\ldots,n-1)
\tag{5}
\end{equation}
gesetzt. Dann ist
\[
 \bar B_i=[a_{i+1}E,B_i]
\]
und daher nach Satz I
\[
 t(A_i\mid B_i)={\psi(B_i)\over\psi(\bar B_i)}\,t(A_i\mid\bar B_i),
\]
ferner ist $[a_{i+1},b_{i+1}]=a_{i+1}$, woraus nach Satz II
\[
 t(A_i\mid\bar B_i)=t(A_{i+1}\mid B_{i+1})
\]
folgt. Berücksichtigt man endlich, dass $t(A_{n-1}\mid B_{n-1})=1$ ist, so gelangt man zu der Formel
\[
\text{III.}\qquad
 t(A\mid B)=
 {\psi(B)\psi(B_1)\cdots\psi(B_{n-2})
 \over
 \psi(\bar B)\psi(\bar B_1)\cdots\psi(\bar B_{n-2})}.
\]
Die Function $v(A\mid B)$ lässt sich unter Anwendung zweier Sätze, welche ein dualistisches Analogon zu den unter I und II angegebenen bilden, ebenfalls als Quotient aus Producten von $\psi$-Functionen darstellen. Dasselbe kann aber auch erreicht werden, indem man die Formeln § 7, I und die soeben für die Function $t(A\mid B)$ gefundene anwendet. Man erhält dann
\[
\text{IV.}\qquad
 v(A\mid B)=
 {\psi(A)\psi(B_1)\cdots\psi(B_{n-2})
 \over
 \psi(\bar B)\psi(\bar B_1)\cdots\psi(\bar B_{n-2})}.
\]
Die weitere Ausrechnung der Functionen $t(A\mid B)$ und $v(A\mid B)$, bei welcher wir uns auf primäre Classen beschränken können, soll erst an späterer Stelle erfolgen.

\section*{§ 9. Darstellung der Function $\bar\gamma$ durch $t$-Functionen, der Function $\bar\delta$ durch $v$-Functionen.}

Nachdem in § 8 gezeigt worden ist, wie sich die Functionen $t$ und $v$ durch $\psi$-Functionen ausdrücken lassen, soll nunmehr dargethan werden, wie man mit Hülfe der Functionen $t$ und $v$ die Functionen $\bar\gamma$ und $\bar\delta$ darstellen kann.

\end{document}
